\section{Day 10: Continuous Density Functions (Oct 08, 2025)}

\begin{definition}[CDF]
    For any random variable $X$, the function $F_X: \mathbb{R} \to [0, 1]$, where $F_X(x) = P(X \leq x)$. 
\end{definition}

\noindent This applies for discrete and absolutely continuous $X$. If $X$ is discrete:
\[
F_X(x) = \sum_{u \leq x} P(X = u) = \sum_{u \leq x} p_X(u)
\]
or if $X$ is absolutely continuous, we have:
\[
F_X(x) = \int_{-\infty}^x f_X(t) \, dt
\]
\begin{simplethm}[`Distribution Function']
    If $f_X$ is continuous, then $F_X' = f_X$
\end{simplethm}

\begin{proof}
    We verify using the fundamental theorem of calculus:
    \[
    F_X'(x) = \frac{d}{dx} \int_{-\infty}^x f_X(t) \, dt = f_X(x)
    \]
\end{proof}

\begin{definition}[Left Continuous]
    A function is left continuous if the left hand limit always exists for all $x$ in its domain. 
\end{definition}

We prefer using $\displaystyle \lim_{ n \to \infty } f\left(a - \frac{1}{n}\right)$ as opposed to $\displaystyle \lim_{ x \to a^- } f(x)$. 

\begin{theorem}[CDF Properties]
    A CDF of a random variable $X$ must have the following properties:
\begin{itemize}
    \item $0 \leq F_X(x) \leq 1$ for all $x \in \mathbb{R}$
    \item $F_X$ is non-decreasing
    \item $\displaystyle \lim_{ x \to -\infty } F_X(x) = 0$
    \item $\displaystyle \lim_{ x \to \infty } F_X(x) = 1$
\end{itemize}
\end{theorem}

\begin{simplethm}
$X$ is a continuous random variable if and only if $F_X$ is a continuous function for all $x$. 
\end{simplethm}

\noindent CDFs of discrete random variables also have nice properties, and we say that they are piecewise-constant.
